%% kapa: exam=1 round=31 year=2020 period=1 form=A subject=civil_law subject_ko=민법 seq=2 number=2
%% source: Q-Net 감정평가사 기출문제 (cst00309) artlSeq=5207219 file=extracted/1008/5207219/A형/A형/2020년도 제31회 감정평가사 1차 시험_1교시_A형.pdf
%% conversion: kordoc 4.24.51 -> tools/gen1.py (auto); review=proofread-sonnet; flags=none
\begin{question}{2}
권리주체에 관한 설명으로 옳지 않은 것은? (다툼이 있으면 판례에 따름)
\begin{choices}
  \item 의사능력은 자신의 행위의 의미와 결과를 합리적으로 판단할 수 있는 정신적 능력으로 구체적인 법률행위와 관련하여 개별적으로 판단되어야 한다.
  \item 어떤 법률행위가 일상적인 의미만으로 알기 어려운 특별한 법률적 의미나 효과를 가진 경우, 이를 이해할 수 있을 때 의사능력이 인정된다.
  \item 현행 민법은 태아의 권리능력에 관하여 일반적 보호주의를 취한다.
  \item 태아의 상태에서는 법정대리인이 있을 수 없고, 법정대리인에 의한 수증행위도 할 수 없다.
  \item 피상속인과 그의 직계비속 또는 형제자매가 동시에 사망한 것으로 추정되는 경우에도 대습상속이 인정된다.
\end{choices}
\end{question}
\answer{3}
